\documentclass[11pt,a4paper]{article}

\usepackage[utf8]{inputenc}
\usepackage[T1]{fontenc}
\usepackage{lmodern}
\usepackage[english]{babel}
\usepackage{geometry}
\usepackage{graphicx}
\usepackage{float}
\usepackage{caption}
\usepackage{amsmath}
\usepackage{amssymb}
\usepackage{booktabs}
\usepackage{hyperref}
\usepackage{subcaption}
\usepackage{xcolor}

\geometry{margin=1in}
\raggedbottom
\emergencystretch=3em

% Figures live in the repo's runs/ directory; this file sits in report/8/,
% so resolve \includegraphics paths relative to the repository root.
\graphicspath{{../../}}

% Graceful fallback for figures produced by runs that may not be pulled yet.
\newcommand{\figorbox}[2]{%
  \IfFileExists{../../#1}{\includegraphics[width=#2\textwidth]{#1}}{%
    \fbox{\parbox[c][3.5cm][c]{#2\textwidth}{\centering\itshape
    \textcolor{gray}{[figure not yet generated --- run the eval and pull it]}}}}}

\title{Empirical Analysis of Transformers on Graph Connectivity \\
       \large Part VIII: Where Is Connectivity Information Combined?}
\author{Edoardo Paccagnella}
\date{\today}

\begin{document}

\maketitle

% =====================================================================
% DRAFT / IN PROGRESS. Section 5 (the unified similarity-read-out battery) is
% WRITTEN with real data (all of it already existed on disk from Report VII's
% own similarity run, sec:res-similarity there -- no new computation needed,
% only a new attention-scores figure that existed as a FILE but was never
% shown in Report VII's rendered text). Section 6 (two cycles) is a STUB:
% code is implemented and smoke-tested, the HPC job is prepared but not yet
% run. Section 7 (outlook) records ideas not yet implemented. See
% istruzioni.md sec 16 for the full brief this report follows.
% =====================================================================

\begin{abstract}
\noindent Report~VII opened the trunk and found a graded, split-specific
completion signal that lets an unbalanced two-chain split resolve far beyond
the architectural distance wall, and showed that a similarity read-out makes
the model's behaviour far more robust to that signal's imperfections than the
linear read-out used everywhere before it. This report does two things.
First, having decided to standardize on the similarity read-out going
forward, it collects the \emph{whole} mechanistic battery --- the behavioural
sweep, the read-out geometry, the attention scores, the node-to-node
contribution, the layer ablations, the falsification test --- in one place
for the first time, rather than split across a linear-read-out-first
narrative. Second, it asks a new question about the mechanism itself: does
the model's apparent reliance on a path's open \emph{endpoints} reflect
something the endpoints \emph{are} (a distinguished degree-1 node), or
something a path merely \emph{has} (an end)? We test this by closing both
segments of the two-chain split into cycles, removing every endpoint while
keeping the split intact.
\end{abstract}

\section{Introduction}
\label{sec:intro}

\subsection{What Reports I--VII established}
\label{sec:recap}

We treat the following as established (documented, multi-seed, in
Reports~I--VII); this report's new work must stay consistent with them.

\begin{itemize}\setlength\itemsep{0.2em}
  \item \textbf{There is a sharp, architectural distance wall at $3^L$} ($=9$
        at $L=2$), which moves with depth and appears even on graph families
        seen $\sim10^8$ times in training (Reports~III--IV).
  \item \textbf{The model computes distance-bounded matrix powering, not a
        bounded traversal}, and a soft, learned ``node budget'' on dense
        structure is a removable data prior, not a second capacity
        (Report~V).
  \item \textbf{Parallel routes and repeated local structure both help,
        within capacity, in data-shaped ways} rather than architecturally for
        free (Report~VI).
  \item \textbf{An unbalanced two-chain split resolves far beyond the
        distance wall; a balanced split does not, with the same weights}
        (Report~VI, Thread~B) --- and opening the trunk shows why
        (Report~VII): on a graph with exactly two components, once one is
        small enough to resolve completely, the model gets an extra,
        graded completion signal for the \emph{other} component, built by
        the same two-layer attention machinery that does ordinary
        beyond-capacity reach, not a separate circuit. The signal survives
        under \emph{both} training distributions tested (disjoint paths,
        pure Erd\H{o}s--R\'enyi) and at three canvas sizes ($n=20,40,64$); a
        label/index shortcut is refuted throughout.
  \item \textbf{The read-out shape decides how much that signal can hurt, not
        whether the trunk has it.} The linear read-out's specific failure
        modes --- a hard reach floor, a cut decision defended only by the
        sign of an unbounded logit --- are properties of \emph{that}
        read-out. The similarity read-out
        ($z_{ij}=\mathrm{scale}\cdot\cos(h_i,h_j)+\mathrm{bias}$) shows the
        same graded, split-specific mechanistic signal (Report~VII,
        \S sec:res-similarity) but converts it into a materially more robust
        behaviour: at $n=40$, reach never drops below $0.91$ (against a
        linear floor of $0.58$--$0.65$) and cut never drops below $0.995$
        (against a linear collapse to $0.15$ for one training distribution
        and to $0.01$--$0.09$ for another) --- a global, fixed cosine margin
        defends the cut decision everywhere a per-pair unbounded logit could
        not.
\end{itemize}

\subsection{The new question: where, and how, is the information combined?}
\label{sec:question}

Every report through Report~VII reads the trunk's computation through a single
lens at a time: how far it reaches, which node budget it respects, which
oracle it resembles, which attention weights or read-out entries move with a
split. What none of them asks directly is the more basic mechanical question:
\emph{at which point in the forward pass does the model already have the
answer} --- has it already decided ``connected'' or ``not'' by the end of
layer~1, or only after layer~2? Does the feed-forward step (rather than
attention) do the actual \emph{combining} of ``I am near this node, and that
node is near the target, therefore\ldots''? This report is framed around that
question: not just \emph{whether} a signal exists, but \emph{where in the
stack} it is assembled. \S\ref{sec:battery} lays the groundwork by putting
every existing mechanistic quantity for the (now-standard) similarity
read-out in one place, including a new \emph{causal} complement to the
existing node-to-node contribution: instead of only asking how sensitive the
final embedding is to a perturbation of an internal representation, it
intervenes directly on the input graph (removing one edge at a time) and
reads off the resulting change (\S\ref{sec:battery-edge}). \S\ref{sec:cycles}
tests one specific structural claim the existing evidence suggests (that a
path's open ends carry special weight) by removing the ends entirely.
\S\ref{sec:outlook} lists the concrete next steps this framing still calls
for --- comparing the raw per-layer embeddings $h^{(1)}, h^{(2)}$ directly
(not only derived quantities) to see whether a two-component ``block''
structure is visible layer by layer, and inspecting the attention output
added to the residual stream directly.

\section{Setup}
\label{sec:setup}

\paragraph{Model.} The same base model as Reports~III--VII: the
RoBERTa-faithful standard transformer, $L=2$, single head,
$d_{\mathrm{model}}=512$, normalised-ReLU attention. \textbf{From this report
on we standardize on the similarity read-out}
($z_{ij}=\mathrm{scale}\cdot\cos(h_i,h_j)+\mathrm{bias}$, a single global
scale and bias, genuinely symmetric in $h_i,h_j$) rather than the linear
read-out used as the default through Report~VII. The justification is
Report~VII's own finding (\S\ref{sec:recap} above): the two-component
completion signal this project has spent two reports characterising is
present under \emph{either} read-out, but the similarity read-out is
categorically more forgiving of the trunk's imperfect, graded cross-component
mixing --- reach does not floor, cut does not collapse, and (for a model that
never saw disconnected graphs in training) the mechanism-vs-behaviour
dissociation is cleaner. Since the read-out is not the thing under
investigation here, we fix it to the one that lets behaviour track the
mechanism most faithfully. We start with the \emph{disjoint-paths} training
distribution (unions of $1$--$4$ paths covering all $n$ nodes, as in every
two-chain result since Report~VI) at $n=40$, four seeds --- the checkpoint set
already used for Report~VII's own similarity comparison. Other training
distributions and canvas sizes are left for a later pass (\S\ref{sec:outlook}).

\paragraph{Notation.} We reuse Report~VII's forward-pass notation without
repeating the full derivation: $h_i^{(0)}$ is node $i$'s read-in embedding
(after the embedding LayerNorm); $h_i^{(1)}, h_i^{(2)}$ its embedding after
each of the two transformer blocks; $S_{\ell,ij}=q_{i}^\top k_{j}/\sqrt{d_h}$
the raw attention score and $\alpha_{\ell,ij}=\tfrac1n\max(0,S_{\ell,ij})$ the
normalised-ReLU attention weight at layer $\ell$; and
$C_{ik}=\|\partial h_i^{(2)}/\partial h_k^{(0)}\|_F$ the exact node-to-node
contribution (real automatic differentiation through the whole forward pass
--- $V$, $W_O$, the residual identity, the feed-forward step, LayerNorm's true
local derivative), introduced in Report~VII \S sec:setup to replace an
attention-only rollout approximation. \S\ref{sec:outlook} discusses a planned
revision of this last quantity.

\paragraph{Test graphs.} Two-chain graphs partitioning all $n=40$ nodes into
two disjoint paths of size $a$ and $n-a$, swept over $a=1,\dots,20$, exactly
as in Report~VII.

\section{The full mechanistic battery, in one setting}
\label{sec:battery}

Report~VII introduced the similarity read-out comparison as one section
(\S sec:res-similarity there) after five sections built around the linear
read-out; consequently several quantities computed for the similarity
checkpoints were never shown together, and one (the attention-score heatmap)
was computed but not rendered in the report text at all. This section
collects the complete battery for the model this project now treats as
standard --- similarity read-out, disjoint-paths training, $n=40$ --- in one
place. Every quantity below reuses Report~VII's already-validated
computation (same scripts, same checkpoints, same four seeds); no new
computation was needed for this section.

\subsection{Behavioural sweep and the raw cosine logit}
\label{sec:battery-sweep}

Table~\ref{tab:r8sweep} and Figure~\ref{fig:r8sweep} give the full split
sweep. Exact match still falls from $1.0$ to $0$ as the split balances, but
reach in the long component never drops below $0.91$ at any split tested and
recovers to $0.995$ by $a=20$ --- there is no floor-and-recovery cliff down to
$0.58$--$0.65$ the way the linear read-out shows at the same canvas. Cut
stays at $\ge0.995$ across the entire sweep. The exact-match values at
$a=8$--$10$ ($0.500$) hide a clean per-seed split rather than a graded
degradation: two of the four seeds solve every graph exactly at each of these
splits, the other two solve none --- a seed-level threshold shift of a few
nodes, not a within-seed partial success.

\begin{table}[H]\centering\small
\begin{tabular}{lccccc}
\toprule
split $(a,\,40{-}a)$ & exact & reach (long) & reach (short) & cut & pred.\ positive rate\\
\midrule
$(1,39)$  & 1.000 & 1.000 & --    & 1.000 & 0.951\\
$(4,36)$  & 1.000 & 1.000 & 1.000 & 1.000 & 0.820\\
$(7,33)$  & 1.000 & 1.000 & 1.000 & 1.000 & 0.711\\
$(8,32)$  & 0.500 & 0.910 & 1.000 & 1.000 & 0.624\\
$(9,31)$  & 0.500 & 0.917 & 1.000 & 1.000 & 0.603\\
$(10,30)$ & 0.500 & 0.925 & 1.000 & 1.000 & 0.584\\
$(14,26)$ & 0.000 & 0.915 & 1.000 & 0.999 & 0.511\\
$(17,23)$ & 0.000 & 0.961 & 1.000 & 0.995 & 0.501\\
$(20,20)$ & 0.000 & 0.995 & 0.995 & 0.996 & 0.499\\
\bottomrule
\end{tabular}
\caption{\emph{Model:} base standard transformer with the \emph{similarity}
read-out ($L=2$, single head, $d_{\mathrm{model}}=512$,
$z_{ij}=\mathrm{scale}\cdot\cos(h_i,h_j)+\mathrm{bias}$), trained on a single
clean distribution of disjoint paths (unions of $1$--$4$ paths covering all
$n$ nodes), four seeds, $n=40$. \emph{Test set:} two-chain graphs split
$(a,\,40{-}a)$, $200$ independently node-relabelled graphs per split per seed.
\emph{Metrics} (mean over seeds): whole-graph exact match; pairwise reach
inside the long/short path (target connected); pairwise accuracy across the
two paths (target disconnected); the model's overall fraction of pairs
predicted connected.}
\label{tab:r8sweep}
\end{table}

\begin{figure}[H]
    \centering
    \figorbox{runs/report7/report7_figs/r7_sweep_and_logit_n40_pathunion_similarity.png}{0.92}
    \caption{\emph{Model/training/test as in Table~\ref{tab:r8sweep}}, full
    split sweep $a=1,\dots,20$. \emph{Top:} behavioural sweep, error bars
    $=$ std across the four seeds. \emph{Bottom:} mean raw $\cos(h_i,h_j)$
    (before scale/bias) for pairs within the long path, within the short
    path, and across the two paths, mean over seeds, error bars $=$ std
    across seeds.}
    \label{fig:r8sweep}
\end{figure}

\subsection{Read-in geometry}
\label{sec:battery-geometry}

There is no $W_{\mathrm{out}}$ for this read-out, only the read-in embedding
$E_{\mathrm{in}}$ and two global scalars, $\mathrm{scale}$ and
$\mathrm{bias}$. Table~\ref{tab:r8wout} and Figure~\ref{fig:r8wout} give its
geometry: row norms close to uniform across the $40$ labels and small
pairwise cosine similarity between different rows (no fixed node-index
shortcut). The decision boundary is $\mathrm{scale}\cdot\cos(h_i,h_j)+
\mathrm{bias}=0$, i.e.\ $\cos(h_i,h_j)=-\mathrm{bias}/\mathrm{scale}\approx0.146$
for this training distribution --- a single, fixed cosine threshold every
pair is measured against.

\begin{table}[H]\centering\small
\begin{tabular}{lcc}
\toprule
 & $\|e_k\|$ (mean, std) & mean $\cos(e_k,e_l)$ ($k\neq l$)\\
\midrule
seed 1000 & $1.89\pm0.09$ & $+0.250$\\
seed 2000 & $1.19\pm0.05$ & $+0.369$\\
seed 3000 & $2.02\pm0.09$ & $+0.255$\\
seed 4000 & $1.36\pm0.18$ & $+0.299$\\
\midrule
mean scale, bias & \multicolumn{2}{c}{$32.6,\ -4.77$}\\
implied threshold $\cos$ & \multicolumn{2}{c}{$\approx0.146$}\\
\bottomrule
\end{tabular}
\caption{\emph{Model:} the four checkpoints of Table~\ref{tab:r8sweep}.
\emph{Test set:} none --- static properties of the trained weights, not
evaluated on any graph. \emph{Metrics:} read-in row norms $\|e_k\|$ (mean
$\pm$ std.\ over the $40$ labels); mean pairwise cosine similarity between
different rows of $E_{\mathrm{in}}$; the two global read-out scalars (mean
over seeds); and the cosine value at which
$\mathrm{scale}\cdot\cos+\mathrm{bias}=0$, the model's fixed connectivity
threshold.}
\label{tab:r8wout}
\end{table}

\begin{figure}[H]
    \centering
    \figorbox{runs/report7/report7_figs/r7_heatmap_weight_geometry_n40_pathunion_similarity.png}{0.85}
    \caption{\emph{Model:} the four checkpoints of Table~\ref{tab:r8sweep}.
    \emph{Test set:} none --- static weights, as in Table~\ref{tab:r8wout}.
    \emph{Metric:} $E_{\mathrm{in}}=W_{\mathrm{in}}^\top$ (raw entries), the
    cosine-similarity matrix $\cos(e_k,e_l)$, and the row norms $\|e_k\|$ by
    label index (one representative seed).}
    \label{fig:r8wout}
\end{figure}

\subsection{Attention scores and normalized-ReLU $\alpha$}
\label{sec:battery-attn}

This figure was computed for the similarity read-out in Report~VII but never
shown there. Figure~\ref{fig:r8attnscores} gives the real attention-score
matrix $S=QK^\top/\sqrt{d_h}$ and the resulting $\alpha$, per layer, at
$a=4$ (solved) and $a=20$ (failed). Layer~$0$ is a narrow band along the
diagonal at both splits --- local, structure-driven, matching every other
condition in this project. Layer~$1$ is where the two splits differ and
where a feature not visible in the linear read-out's version of this figure
stands out clearly: at \emph{both} $a=4$ and $a=20$, the columns nearest the
two path endpoints (the first few and last few base-node indices) draw a
markedly higher $\alpha$ than the rest of the row from almost every other
node --- bright vertical bands at both edges of the layer-1 panels --- while
the component boundary itself is visible as a break in the near-diagonal band
only once the short component is a large enough fraction of the canvas
(clearest at $a=20$).

\begin{figure}[H]
    \centering
    \figorbox{runs/report7/report7_figs/r7_heatmap_attention_scores_n40_pathunion_similarity.png}{0.95}
    \caption{\emph{Model:} one seed of the checkpoints of Table~\ref{tab:r8sweep}.
    \emph{Test set:} one representative two-chain graph at $a=4$ and one at
    $a=20$, node order fixed to base graph position (short path first, then
    long) for readability. \emph{Metric:} the real attention score matrix
    $S=QK^\top/\sqrt{d_h}$ and the resulting normalized-ReLU attention
    $\alpha=\mathrm{ReLU}(S)/n$, layer $0$ (top) and layer $1$ (bottom),
    averaged over $80$ independent random relabellings of the same
    underlying graph then mapped back to base node order.}
    \label{fig:r8attnscores}
\end{figure}

\subsection{Node-to-node contribution: internal sensitivity, not yet a causal claim}
\label{sec:battery-contrib}

Figure~\ref{fig:r8leak} gives the leak fraction (the share of a long-path
node's exact contribution mass $C_{ik}$ that traces back to the short
component) across the full sweep, and Figure~\ref{fig:r8contribmat} the raw
matrix at the same two representative splits. The leak fraction jumps from
$0\%$ at $a=1$ to a \emph{flat} plateau of $\approx21$--$23\%$ from $a=4$
onward --- neither the smooth monotone rise nor the sharp spike the linear
read-out shows at the same canvas (Report~VII \S sec:res-attention). The
matrix itself shows a genuine, non-trivial band of cross-component mixing at
\emph{both} splits, even though behaviour is nearly perfect at $a=4$: a real
dissociation between mechanism and behaviour --- the trunk mixes across the
component boundary by a similar amount whether or not the split is solved,
and it is the read-out, not the mixing, that decides whether that costs
anything (\S\ref{sec:recap}). One caveat about what this quantity actually
shows, worth stating plainly rather than only in a footnote: $C_{ik}$ is the
sensitivity of $h_i^{(2)}$ to an \emph{arbitrary} perturbation of $h_k^{(0)}$,
and $h_k^{(0)}$ is itself an internal representation built from node $k$'s
\emph{whole} neighbourhood (through the read-in and a LayerNorm), not a
quantity tied only to node $k$'s own identity. A large $C_{ik}$ is therefore
good evidence of internal token-to-token \emph{sensitivity}, but not by
itself evidence that node $k$'s actual incident edges have a \emph{causal}
effect on node $i$ --- that question is taken up directly in
\S\ref{sec:battery-edge}, which intervenes on the input graph itself rather
than on an internal embedding.

\begin{figure}[H]
    \centering
    \figorbox{runs/report7/report7_figs/r7_attention_leak_n40_pathunion_similarity.png}{0.7}
    \caption{\emph{Model/training as in Table~\ref{tab:r8sweep}}. \emph{Test
    set:} two-chain graphs at the splits shown, $64$ independently
    node-relabelled graphs per seed. \emph{Metric:} for each long-component
    node, the fraction of its exact node-to-node contribution mass
    $C_{ik}=\|\partial h_i^{(2)}/\partial h_k^{(0)}\|_F$ (\S\ref{sec:setup})
    that traces back to the short component rather than its own long
    component, mean over long-component nodes; small dots are individual
    seeds, the line is the mean over seeds.}
    \label{fig:r8leak}
\end{figure}

\begin{figure}[H]
    \centering
    \figorbox{runs/report7/report7_figs/r7_heatmap_contrib_matrix_n40_pathunion_similarity.png}{0.95}
    \caption{\emph{Model:} one seed of the checkpoints of Table~\ref{tab:r8sweep}
    (representative; the pattern is consistent across seeds).
    \emph{Test set:} two-chain graphs at the split shown ($a=4$ left, $a=20$
    right), $64$ independently node-relabelled graphs, $C_{ik}$ averaged
    over them. \emph{Metric:} the raw exact node-to-node contribution matrix
    $C_{ik}=\|\partial h_i^{(2)}/\partial h_k^{(0)}\|_F$, rows/columns
    ordered by base position (short component first, then long). Dotted
    white lines mark the component boundary.}
    \label{fig:r8contribmat}
\end{figure}

\subsection{A causal complement: edge-ablation contribution}
\label{sec:battery-edge}

\paragraph{Why a second measure.} \S\ref{sec:battery-contrib} measures
sensitivity to a perturbation of an already-processed embedding. To make a
\emph{causal} statement about the short component's edges --- do they
actually drive the long component's beyond-capacity predictions, or does the
trunk merely correlate the two --- the natural experiment intervenes on the
input graph directly: remove one edge, re-run the forward pass, and measure
what changes.

\paragraph{Construction.} For every edge $e=\{u,v\}$ of a two-chain graph
(both directions removed from $A'$ at once, since the model is trained on
symmetric inputs; self-loops on the diagonal are left untouched), we compare
a forward pass on the intact graph against one with $e$ removed, and read
three quantities off node $i$: the change in its final embedding
$\|h_i^{(2)}(A'\!\setminus\!e)-h_i^{(2)}(A')\|_2$; the mean absolute change
over its whole logit row; and the same logit change restricted to the pairs
$(i,j)$ that matter most for the beyond-capacity puzzle --- both $i,j$ inside
the long component, at true graph distance $>9$. Averaging (never summing,
so that a path endpoint's degree of $1$ is not confounded with an internal
node's degree of $2$) over every edge touching a given node $k$ turns each of
the three quantities into a $40\times40$ matrix, column $k$ = which node's
incident edges were removed, row $i$ = whose state we observe --- the same
shape as Figure~\ref{fig:r8contribmat}, but built from real interventions on
the graph rather than a derivative of an internal embedding. One honest
caveat carried through the design rather than left implicit: on a path or a
cycle \emph{every} edge is a bridge for at least one pair, so removing it
also changes that pair's true connectivity target --- this measures the
effect of a real structural change, not an independent, additive
contribution the way the phrase ``edge contribution'' might suggest.

\emph{[Data pending --- the computation is implemented and smoke-tested (a
forward-pass-only intervention, no backward pass, cheaper than
\S\ref{sec:battery-contrib}'s Jacobian) but not yet run on the real
checkpoints at the time of writing. Replace this paragraph with the three
heatmaps at $a=4$ and $a=20$, the edge/logit/far leak-fraction curve across
the full split sweep, and a comparison against
\S\ref{sec:battery-contrib}'s reading once the job returns. The prediction to
check: at $a=4$ the short component's edges should show a real effect on the
long component's beyond-capacity logits (a high $C^{\mathrm{far}}$ leak); at
$a=20$ that effect should be small and confined near the diagonal.]}

\subsection{Layer ablation}
\label{sec:battery-ablation}

Table~\ref{tab:r8ablation} and Figure~\ref{fig:r8ablation} give the ablation
sweep. The bypass condition (both transformer blocks skipped, read-out
applied directly to the read-in embedding) already gives cut $0.95$--$0.995$
across the whole sweep --- higher than several of the attention-ablated
conditions --- while its reach is low throughout ($0.10$--$0.20$): most of
what makes cut work is already present in the read-in embeddings' cosine
geometry before any attention runs, and the trunk's main job for this
read-out is building reach, not cut. Zeroing attention layer~$1$ knocks reach
down only mildly ($0.82$--$0.89$) while zeroing layer~$0$ leaves reach
essentially untouched ($0.987$--$0.992$) but collapses cut to $0.01$--$0.04$
once the split is no longer tiny --- the reverse of the dependency the
linear read-out showed, where cut depended on layer-$1$ feed-forward
specifically.

\begin{table}[H]\centering\small
\begin{tabular}{lcccccccc}
\toprule
condition & $a{=}1$ & $a{=}4$ & $a{=}7$ & $a{=}8$ & $a{=}10$ & $a{=}14$ & $a{=}17$ & $a{=}20$\\
\midrule
\multicolumn{9}{l}{\emph{reach (long component)}}\\
baseline    & 1.000 & 1.000 & 1.000 & 0.911 & 0.925 & 0.915 & 0.961 & 0.995\\
zero attn 0 & 0.992 & 0.992 & 0.992 & 0.992 & 0.991 & 0.990 & 0.988 & 0.987\\
zero attn 1 & 0.885 & 0.876 & 0.868 & 0.865 & 0.859 & 0.844 & 0.831 & 0.820\\
zero mlp 0  & 0.959 & 0.954 & 0.953 & 0.959 & 0.949 & 0.942 & 0.937 & 0.923\\
zero mlp 1  & 0.794 & 0.944 & 0.913 & 0.886 & 0.902 & 0.921 & 0.961 & 0.991\\
bypass      & 0.102 & 0.110 & 0.120 & 0.124 & 0.132 & 0.152 & 0.172 & 0.197\\
\midrule
\multicolumn{9}{l}{\emph{cut (target: disconnected)}}\\
baseline    & 1.000 & 1.000 & 1.000 & 1.000 & 1.000 & 1.000 & 0.996 & 0.997\\
zero attn 0 & 0.997 & 0.035 & 0.025 & 0.024 & 0.019 & 0.015 & 0.013 & 0.014\\
zero attn 1 & 0.752 & 0.903 & 0.382 & 0.343 & 0.293 & 0.241 & 0.225 & 0.219\\
zero mlp 0  & 0.917 & 0.047 & 0.136 & 0.113 & 0.100 & 0.082 & 0.074 & 0.074\\
zero mlp 1  & 0.807 & 0.226 & 0.335 & 0.477 & 0.518 & 0.686 & 0.739 & 0.760\\
bypass      & 0.949 & 0.986 & 0.991 & 0.992 & 0.993 & 0.994 & 0.995 & 0.995\\
\bottomrule
\end{tabular}
\caption{\emph{Model/training as in Table~\ref{tab:r8sweep}}, four seeds.
\emph{Test set:} two-chain graphs at the splits shown, $150$ graphs per seed,
independently node-relabelled. \emph{Metric:} reach (long component, top
block) and cut (bottom block), mean over seeds, under six forward-pass
conditions --- unmodified (baseline); the attention or feed-forward branch
zeroed at layer $0$ or $1$; and both transformer blocks skipped entirely
(bypass, read-out applied directly to the read-in embedding).}
\label{tab:r8ablation}
\end{table}

\begin{figure}[H]
    \centering
    \figorbox{runs/report7/report7_figs/r7_ablation_reach_cut_n40_pathunion_similarity.png}{0.95}
    \caption{Table~\ref{tab:r8ablation} as curves. \emph{Model/training/test
    as in Table~\ref{tab:r8ablation}}, mean over four seeds. Red dashed
    line: capacity $3^L=9$.}
    \label{fig:r8ablation}
\end{figure}

\subsection{Falsification test: three components}
\label{sec:battery-threeway}

Table~\ref{tab:r8threeway} repeats Report~VII's falsification construction: a
small, cleanly-resolvable path plus two large paths that are \emph{not}
connected to each other. If the model implemented a crude ``not the small one
therefore connected'' rule, it would wrongly merge the two large components.
It does not, at any small-component size tested: $\mathrm{cut}(\mathrm{L}_1,
\mathrm{L}_2)\ge0.988$ throughout, reaching $1.000$ for several cells, with
reach inside each large component close to perfect once the small component
is not itself a single node.

\begin{table}[H]\centering\small
\begin{tabular}{lccc}
\toprule
small & reach L$_1$/L$_2$ & cut(S,L) & cut(L$_1$,L$_2$)\\
\midrule
1  & 0.998 & 1.000 & 0.998\\
2  & 1.000 & 1.000 & 0.988\\
4  & 1.000 & 1.000 & 0.999\\
7  & 1.000 & 1.000 & 0.999\\
8  & 1.000 & 1.000 & 1.000\\
10 & 1.000 & 1.000 & 1.000\\
\bottomrule
\end{tabular}
\caption{\emph{Model/training as in Table~\ref{tab:r8sweep}}. \emph{Test
set:} three disjoint paths partitioning all $40$ nodes --- one small path
(size shown) and two large paths, the two large paths \emph{not} connected
to each other; $300$ graphs per seed, independently node-relabelled.
\emph{Metrics} (mean over four seeds): pairwise reach inside the two large
components (mean of L$_1$, L$_2$, target connected); pairwise accuracy
between the small path and either large path (target disconnected, mean);
pairwise accuracy between the two large paths (target disconnected --- the
decisive column).}
\label{tab:r8threeway}
\end{table}

\section{Testing the path-endpoint hypothesis: closing the chains into cycles}
\label{sec:cycles}

\paragraph{Where this comes from.} Two independent pieces of evidence, both
in \S\ref{sec:battery} and in Report~VII, single out a path's open
\emph{endpoints} (its two degree-$1$ nodes) as structurally unusual. Under the
linear read-out (Report~VII \S sec:res-heatmaps), the far endpoint of each
path segment is an outsized \emph{source} of node-to-node contribution ---
$3$--$4\times$ the per-node average --- without being an unusual
\emph{receiver}, read there as a ``this path is closed'' signal broadcast
outward. Under the similarity read-out (\S\ref{sec:battery-attn} above), the
columns nearest each path's two endpoints draw markedly more layer-$1$
attention than the rest of the graph, at every split tested. Both readings
point the same way, but neither rules out the alternative: that these effects
are an artefact of \emph{how the test graphs are built} (a path always has
exactly two special nodes) rather than something the model's completion
mechanism actually depends on.

\paragraph{The test.} We remove every endpoint while changing nothing else:
close each of the two path segments into a \textbf{cycle}, keeping the same
split $(a,\,n-a)$, the same total node count, and the same
disjoint-two-component structure --- only now neither component has a
degree-$1$ node. % generate_split_cycles_graph(n, short_len) in data.py, the
                  % direct structural analogue of generate_split_chains_graph
                  % used for every two-chain result since Report VI.
A cycle needs at least three nodes, so the sweep starts at
$a=3$ rather than $a=1$. We feed this construction to the \emph{same}
checkpoints as \S\ref{sec:battery} (disjoint-paths-trained, similarity
read-out, $n=40$) --- eval-only, no retraining --- and read the identical
battery: the behavioural sweep, the exact-Jacobian contribution leak
fraction of \S\ref{sec:battery-contrib}, and the causal edge-ablation
contribution of \S\ref{sec:battery-edge}, all with no change to either
measure's definition --- only the underlying graph topology changes. One
adjustment the cycle geometry itself forces: a $k$-cycle's diameter is
$\lfloor k/2\rfloor$, roughly half a $k$-node path's, so the
beyond-capacity pairs the far-logit variant of \S\ref{sec:battery-edge}
looks at only exist once the long \emph{cycle} segment is long enough
(around $19$ nodes, versus around $10$ for a path of the same role).

\paragraph{Two readings, stated before looking at the data.} If the model
solves the cycle sweep about as well as the chain sweep (same split
threshold, similar reach/cut pattern), the endpoints were \emph{incidental}:
whatever the completion mechanism actually tracks does not depend on a
degree-$1$ node existing, and the source/high-attention effect at the
endpoints was a by-product of them being the two nodes farthest from every
other node on their segment, not a special role tied to being an open end.
If instead the model's performance degrades sharply or the split threshold
moves, the endpoints are \emph{load-bearing}: removing them removes
something the mechanism relies on to decide a segment is fully resolved.

\emph{[Data pending --- the eval is prepared (eval-only on the existing
similarity checkpoints, no training) but not yet run at the time of writing.
Replace this paragraph with the real sweep table, the leak-fraction
comparison against Figure~\ref{fig:r8leak}, the edge-ablation comparison
against \S\ref{sec:battery-edge}, and the verdict once the job returns.]}

\section{Outlook}
\label{sec:outlook}

The causal edge-ablation contribution of \S\ref{sec:battery-edge} answers the
concern that motivated it (\S\ref{sec:battery-contrib}): it is implemented and
its design is fixed, only the data collection is still pending. Beyond it,
three concrete next steps follow from the framing of \S\ref{sec:question},
none implemented yet.

\begin{itemize}\setlength\itemsep{0.15em}
  \item \textbf{A chord-addition control for the edge-ablation measure.} On a
        path or a cycle every edge is a bridge, so removing one always changes
        the true connectivity target for at least one pair
        (\S\ref{sec:battery-edge}) --- the measured effect is real but
        confounded with that target change. Repeating the same intervention
        with an edge \emph{added} between two nodes already in the same
        component (a chord) leaves the true target unchanged and would isolate
        the effect of a structural change from the effect of a target change.
  \item \textbf{Plot the per-layer embeddings themselves, not only derived
        quantities.} Every mechanistic result so far (attention, contribution,
        read-out logits) is a quantity \emph{derived from} $h^{(1)}$ and
        $h^{(2)}$. Plotting the embeddings directly (e.g.\ a low-dimensional
        projection, coloured by which component a node belongs to) at each
        layer would show directly whether a two-component ``block''
        structure is visible in the representation itself, and at which
        layer it first appears.
  \item \textbf{Inspect where the information is combined.} The working
        hypothesis is that the trunk first establishes, per node, something
        like ``which other nodes are near me'' and only later \emph{combines}
        that with the target's identity to decide connectedness --- and that
        this combination step is a natural place to look for in the
        feed-forward sub-layer rather than attention. Plotting the attention
        output added to the residual stream directly (rather than only the
        attention weights $\alpha$) at each layer is the first concrete probe
        of this, to see whether the two branches are already combined before
        the feed-forward step or only after it.
\end{itemize}

Beyond these three, the battery of \S\ref{sec:battery} and the cycle test of
\S\ref{sec:cycles} are both, for now, run only on the disjoint-paths
training distribution at $n=40$. Extending either to the pure-Erd\H{o}s--R\'enyi
distribution and to $n=20/64$ (as Report~VII did for the linear read-out) is
the natural next pass once the questions above are settled.

\end{document}
